% 调和函数（综述）
% license CCBYSA3
% type Wiki

本文根据 CC-BY-SA 协议转载翻译自维基百科\href{https://en.wikipedia.org/wiki/Harmonic_function}{相关文章}。


在数学、数理物理以及随机过程理论中，调和函数（harmonic function）是指满足拉普拉斯方程（Laplace's equation）的函数。  
设
\[
f : U \to \mathbb{R}~
\]
其中 \(U \subseteq \mathbb{R}^n\) 为一个开集，且 \(f\) 是定义在 \(U\) 上的二阶连续可微函数（即 \(f \in C^2(U)\)）。若 \(f\) 在 \(U\) 上处处满足
\[
\frac{\partial^2 f}{\partial x_1^2}
+ \frac{\partial^2 f}{\partial x_2^2}
+ \cdots
+ \frac{\partial^2 f}{\partial x_n^2}
= 0~
\]
则称 \(f\) 为一个\textbf{调和函数}。  
该方程通常记作
\[
\nabla^2 f = 0,
\quad \text{或等价地} \quad
\Delta f = 0.~
\]
\subsection{“Harmonic” 一词的词源}
“调和函数（harmonic function）”中的形容词“harmonic”（意为“谐和的”或“调和的”）起源于一根被拉紧的弦上某一点所经历的谐振运动（harmonic motion）。描述此类运动的微分方程的解可以用正弦函数（sine）与余弦函数（cosine）来表示，因此这些函数被称为谐波（harmonics）。傅里叶分析（Fourier analysis）研究如何将单位圆上的函数展开为这些谐波的级数。若考虑高维空间中单位 \(n\) 球面上的谐波类比，便得到所谓的球谐函数（spherical harmonics）。这些函数满足拉普拉斯方程。随着时间的推移，“harmonic”一词被推广，用来指代所有满足拉普拉斯方程的函数。\(^\text{[1]}\)  
\subsection{例子}
以下给出了一些二元调和函数的典型例子：
\begin{enumerate}
  \item \textbf{任意全纯函数（holomorphic function）的实部或虚部。}  
  实际上，定义在平面上的所有调和函数都可以表示为某个全纯函数的实部或虚部。
  \item 函数$f(x, y) = e^{x} \sin y$.这是前述情况的一个特例，因为$f(x, y) = \operatorname{Im}\!\left(e^{x + i y}\right)$,而 \( e^{x + i y} \) 是一个全纯函数。对 \(x\) 的二阶偏导数为$\frac{\partial^2 f}{\partial x^2} = e^{x} \sin y$,对 \(y\) 的二阶偏导数为$\frac{\partial^2 f}{\partial y^2} = -\, e^{x} \sin y$,因此$\frac{\partial^2 f}{\partial x^2} + \frac{\partial^2 f}{\partial y^2} = 0$,故 \(f\) 是调和的。
  \item 函数$f(x, y) = \ln\!\left(x^2 + y^2\right),\quad (x, y) \in \mathbb{R}^2 \setminus \{0\}$.该函数定义在去掉原点的平面上，可用于描述：一条线电荷（line charge）产生的电势；或一根无限长圆柱体质量（cylindrical mass）产生的引力势。
\end{enumerate}
对于三元调和函数，常见的例子列于下表。设$r^2 = x^2 + y^2 + z^2$,则可构造出许多满足拉普拉斯方程$\Delta f = 0$的函数，例如球对称势函数与球谐函数。

\begin{table}[ht]
\centering
\caption{}\label{tab_THhs1}
\begin{tabular}{|c|c|}
\hline
\textbf{函数}&\textbf{奇点}\\
\hline
\(\displaystyle \frac{1}{r}\)& 位于原点的单位点电荷 \\
\hline
\(\displaystyle \frac{x}{r^{3}}\)& 位于原点的沿 \(x\) 方向偶极子 \\
\hline
\(\displaystyle -\ln\!\left(r^{2}-z^{2}\right)\) & 整个 \(z\) 轴上的单位线电荷 \\
\hline
\(\displaystyle -\ln(r+z)\) & 负 \(z\) 轴上的单位线电荷 \\
\hline
\(\displaystyle \frac{x}{r^{2}-z^{2}}\) & 整个 \(z\) 轴上的沿 \(x\) 方向偶极子线 \\
\hline
\(\displaystyle \frac{x}{r(r+z)}\) & 负 \(z\) 轴上的沿 \(x\) 方向偶极子线 \\
\hline
\end{tabular}
\end{table}
在物理学中出现的调和函数（harmonic functions）由其奇点（singularities）和边界条件（boundary conditions）决定，例如狄利克雷边界条件（Dirichlet boundary conditions）或诺伊曼边界条件（Neumann boundary conditions）。在无边界区域上，如果给定一个调和函数，可以通过在其上加上任意整函数（entire function）的实部或虚部，得到具有相同奇点的另一个调和函数。因此，在这种情况下，调和函数并不由其奇点唯一确定；然而，在物理情形中，我们通常要求当\( r \to \infty \) 时，解趋于 \(0\)。在这种情况下，解的唯一性由刘维尔定理（Liouville's theorem）保证。

上述调和函数的奇点可以用静电学（electrostatics）的术语描述为“电荷”（charges）和“电荷密度”（charge densities），因此相应的调和函数与这些电荷分布产生的静电势（electrostatic potential）成比例。若将任意上述函数乘以常数、旋转，或加上常数，所得仍是一个调和函数。对这些函数作反演（inversion），将得到另一个调和函数，其奇点为原奇点关于某个球面“镜像”的位置。此外，任意两个调和函数的和仍然是一个调和函数。

最后，若考虑 \(n\) 个变量的情形，调和函数的例子包括：
\begin{enumerate}
\item 在整个 \(\mathbb{R}^{n}\) 上的常数函数、线性函数与仿射函数（affine functions），例如电容器极板间的电势，或平板质量分布产生的引力势；
\item 定义在 \(\mathbb{R}^{n}\setminus\{0\}\) 上的函数$f(x_{1}, \dots, x_{n}) = \left(x_{1}^{2} + \cdots + x_{n}^{2}\right)^{1 - n/2}, \quad n > 2$,它也是一个调和函数。
\end{enumerate}
\subsection{性质}
在给定开集 \( U \) 上的所有调和函数（harmonic functions）的集合，可以看作拉普拉斯算子 \(\Delta\) 的核，因此它构成一个实数域 \(\mathbb{R}\) 上的向量空间：任意调和函数的线性组合仍然是调和函数。

若 \( f \) 是定义在 \( U \) 上的调和函数，则 \( f \) 的所有偏导数在 \( U \) 上仍是调和函数。拉普拉斯算子 \(\Delta\) 与偏导算子在此类函数上可交换。

在许多方面，调和函数是解析函数（holomorphic functions）的实数对应物。所有调和函数都是解析的（analytic），即它们在局部可表示为幂级数展开。这一事实是椭圆算子（elliptic operators）的一般性质，而拉普拉斯算子是其中最典型的例子。

收敛的调和函数序列的逐点一致极限仍然是调和函数。这是因为任意满足平均值性质（mean value property）的连续函数都是调和的。例如，在区域 \((-\infty, 0) \times \mathbb{R}\) 上定义序列$f_{n}(x,y) = \frac{1}{n} e^{nx} \cos(ny)$,该序列中的每个函数都是调和的，并一致收敛于零函数。但注意其偏导数并不一致收敛于零函数（即零函数的导数）。此例说明在判断极限函数是否调和时，依赖平均值性质与连续性的重要性。
\subsection{与复变函数论的联系}
任意全纯函数（holomorphic function）的实部与虚部都是 \(\mathbb{R}^2\) 上的调和函数（称为一对共轭调和函数，harmonic conjugates）。反之，任意定义在 \(\mathbb{R}^2\) 的开子集 \(\Omega\) 上的调和函数 \(u\)，在局部都是某个全纯函数的实部。事实上，若写作 \( z = x + i y \)，则定义复函数\(g(z) := u_x - i u_y\),因为它满足柯西–黎曼方程（Cauchy–Riemann equations），所以 \( g \) 在 \(\Omega\) 内是全纯的。因此，\( g \) 在局部有一个原函数 \( f \)，并且 \( f' = g \)，从而 \( u \) 是 \( f \) 的实部（差一个常数）。

尽管上述与全纯函数的一一对应关系仅在两个实变量的情形成立，但 \( n \) 个变量的调和函数仍保留了许多典型的解析函数性质。例如：  它们是实解析的（real analytic）；满足最大值原理（maximum principle）与平均值原理（mean-value principle）；存在奇点可去定理（theorem of removal of singularities）与刘维尔定理（Liouville theorem），其形式与复变函数论中的相应定理完全类似。
\subsection{调和函数的性质}
许多调和函数的重要性质可以直接从拉普拉斯方程（Laplace's equation）推导出来。
\subsubsection{调和函数的正则性定理}
调和函数在开集上是无限可微的（infinitely differentiable）。实际上，调和函数是实解析函数（real analytic）的，也就是说，它们在局部可以表示为收敛幂级数。
\subsubsection{最大值原理}
调和函数满足以下最大值原理：若 \( K \subset U \) 是非空紧集（compact subset），则函数 \( f \) 在 \( K \) 上的限制 \( f|_{K} \) 的最大值与最小值都在 \( K \) 的边界上取得。若 \( U \) 是连通的（connected），则这意味着：调和函数 \( f \) 不可能在内部点处取得局部极大值或局部极小值，除非 \( f \) 是常数函数。类似的性质也可以推广到次调和函数（subharmonic functions）。
\subsubsection{平均值性质}
若 \( B(x, r) \) 是一个以 \( x \) 为中心、半径为 \( r \) 的球，且该球完全包含在开集 \(\Omega \subset \mathbb{R}^n\) 中，则调和函数\(u:\Omega \to \mathbb{R}\)在球心处的值 \(u(x)\)，等于函数 \(u\) 在球面上的平均值；该平均值也等于 \(u\) 在球体内部的平均值。换言之：
\[
u(x)
=
\frac{1}{n\omega_n r^{n-1}}
\int_{\partial B(x,r)} u\,d\sigma
=
\frac{1}{\omega_n r^n}
\int_{B(x,r)} u\,dV~
\]
其中 \(\omega_n\) 是 \(n\) 维单位球的体积，\(\sigma\) 表示 \((n-1)\) 维的面积测度。

反之，任何满足（体积）平均值性质的局部可积函数（locally integrable function）都是无限可微的并且是调和的。

卷积形式（Formulation in Terms of Convolutions）设
\[
\chi_r := 
\frac{1}{|B(0,r)|}\chi_{B(0,r)}
=
\frac{n}{\omega_n r^n}\chi_{B(0,r)}~
\]
其中 \(\chi_{B(0,r)}\) 是以原点为中心、半径为 \(r\) 的球的特征函数（characteristic function），并且归一化使得\(\int_{\mathbb{R}^n}\chi_r\,dx = 1\).则函数 \(u\) 在 \(\Omega\) 上是调和的，当且仅当对所有满足 \(B(x,r) \subset \Omega\) 的 \(x, r\)，有
\[
u(x) = (u * \chi_r)(x).~
\]

\textbf{证明思路（Sketch of the Proof）}调和函数的平均值性质及其逆命题可以从下式直接推出：对任意 \(0 < s < r\)，考虑非齐次方程
\[
\Delta w = \chi_r - \chi_s.~
\]
该方程存在一个显式解 \(w_{r,s}\)，它属于 \(C^{1,1}\) 类并具有紧支集于 \(B(0,r)\)。

若 \(u\) 在 \(\Omega\) 上调和，则有
\[
0 = \Delta u * w_{r,s} = u * \Delta w_{r,s} = u * \chi_r - u * \chi_s~
\]
在集合
\[
\Omega_r := \{x \in \Omega \mid \operatorname{dist}(x, \partial \Omega) > r \}~
\]
中成立。

由于 \(u\) 在 \(\Omega\) 上连续，且当 \(s \to 0\) 时 \(u * \chi_s \to u\)，这说明 \(u\) 满足平均值性质。反之，若 \(u \in L^1_{\mathrm{loc}}(\Omega)\) 且满足平均值性质，即
\[
u * \chi_r = u * \chi_s~
\]
在所有 \(0 < s < r\) 的情形下于 \(\Omega_r\) 成立，则反复卷积 \(m\) 次可得
\[
u = u * \chi_r = u * \chi_r * \cdots * \chi_r, \qquad x \in \Omega_{mr}~
\]
从而 \(u \in C^{m-1}(\Omega_{mr})\)，因为 \(\chi_r\) 的 \(m\)-重卷积属于 \(C^{m-1}\) 类并支撑于 \(B(0,mr)\)。由于 \(r,m\) 可任取，得 \(u \in C^{\infty}(\Omega)\)。

此外，
\[
\Delta u * w_{r,s} = u * \Delta w_{r,s} = u * \chi_r - u * \chi_s = 0~
\]
对所有 \(0 < s < r\) 成立，因此由变分法基本定理可得 \(\Delta u = 0\)，即 \(u\) 是调和函数。这证明了调和性与平均值性质的等价性。

推广形式（Generalized Form）若 \(h\) 是一个支撑在 \(B(x, r)\) 中的球对称函数（spherically symmetric function），且\(\int h = 1\),则有\(u(x) = (h * u)(x)\).换句话说，我们可以取函数 \(u\) 在某点周围的加权平均值，并得到 \(u(x)\) 本身。特别地，若取 \(h\) 为 \(C^{\infty}\) 函数，则即使我们只知道 \(u\) 在分布意义下的行为，也可以恢复其点值。参见魏尔引理（Weyl’s lemma）。
\subsubsection{哈纳克不等式}
设
\[
V \subset \overline{V} \subset \Omega~
\]
为有界区域 \(\Omega\) 内的连通集。  
则对于任意非负的调和函数 \(u\)，存在仅依赖于 \(V\) 和 \(\Omega\) 的常数 \(C > 0\)，使得哈纳克不等式成立：
\[
\sup_{V} u \leq C \inf_{V} u.~
\]
此结果刻画了调和函数在连通区域内的值变化受控性：  
在有限区域中，若 \(u\) 不为零且调和，则它在整个区域内不会有剧烈波动。
\subsubsection{奇点可去原理}
对于调和函数，以下奇点可去原理成立：若 \( f \) 是定义在穿孔开集\(\Omega \setminus \{x_0\} \subset \mathbb{R}^n\)上的调和函数，且在点 \(x_0\) 处的奇性比基本解（fundamental solution）更弱（当 \(n > 2\) 时），即满足
\[
f(x) = o\bigl(|x - x_0|^{2 - n}\bigr), \qquad \text{当 } x \to x_0~
\]
则 \(f\) 可以扩展为定义在整个 \(\Omega\) 上的调和函数。此结果与复分析中黎曼定理（Riemann's theorem）关于复函数奇点可去性的结论相对应：若一个全纯函数在孤立点的奇性足够“温和”，则该奇点是可去的。
\subsubsection{刘维尔定理}
\textbf{定理：}若 \( f \) 是定义在整个 \(\mathbb{R}^n\) 上的调和函数，且 \( f \) 有上界或有下界，则 \( f \) 必为常数函数。

（可与复变函数中的刘维尔定理（Liouville's theorem for holomorphic functions）相比较。）

Edward Nelson 的简短证明（Nelson’s Short Proof）Edward Nelson 给出了该定理在“有界调和函数”情况下的一个极为简短的证明\(^\text{[2]}\)，该证明利用了前述的平均值性质。

取任意两点，在这两点分别作半径相同的两个球。若半径取足够大，则这两个球除极小部分外几乎重合。由于 \( f \) 是有界函数，因此其在这两个球上的平均值任意接近。  由平均值性质可知，\( f \) 在任意两点处取相同值。于是 \( f \) 必为常数。

推广到“仅有上界或下界”的情形

若调和函数 \( f \) 仅有上界或下界，证明可作如下调整。通过平移与取相反数，我们可假设 \( f \ge 0 \)。设 \(x, y \in \mathbb{R}^n\)，且令任意 \(R > 0\)，定义\(r = R + d(x, y)\),其中 \(d(x,y)\) 表示两点间的距离。由三角不等式可知：\(B_R(x) \subset B_r(y)\).

利用平均值性质与积分的单调性，有
\[
f(x)
=
\frac{1}{\operatorname{vol}(B_R)}
\int_{B_R(x)} f(z)\,dz
\le
\frac{1}{\operatorname{vol}(B_R)}
\int_{B_r(y)} f(z)\,dz.~
\]
注意到 \(\operatorname{vol}(B_R)\) 与 \(x\) 无关，因此可记为常量。进一步，将分子与分母同时乘除以 \(\operatorname{vol}(B_r)\)，再用平均值性质，可得：
\[
f(x)
\le
\frac{\operatorname{vol}(B_r)}{\operatorname{vol}(B_R)} f(y).~
\]
而当 \(R \to \infty\) 时，
\[
\frac{\operatorname{vol}(B_r)}{\operatorname{vol}(B_R)}
=
\frac{(R + d(x,y))^n}{R^n}
\to 1.~
\]
因此得到\(f(x) \le f(y)\).同理，交换 \(x, y\) 的角色可得\(f(y) \le f(x)\),于是\(f(x) = f(y)\).故 \( f \) 在整个 \(\mathbb{R}^n\) 上恒为常数。

概率论证明（Probabilistic Proof）另一种证明方法利用布朗运动（Brownian motion）：设 \(\{B_t\}_{t \ge 0}\) 是在 \(\mathbb{R}^n\) 上的布朗运动，且\(B_0 = x_0\).则对任意 \(t \ge 0\)，有\(\mathbf{E}[f(B_t)] = f(x_0)\).这表明调和函数 \(f\) 对应于布朗运动的一个鞅（martingale）。利用概率耦合（probabilistic coupling）论证即可得出相同结论：若 \(f\) 有界（或有上界/下界），则其为常数。\(^\text{[3]}\)
\subsection{推广形式}
\subsubsection{弱调和函数}
若函数（或更一般地，分布）在弱意义下满足拉普拉斯方程
\[
\Delta f = 0~
\]
则称其为弱调和函数（weakly harmonic function）。即：\(f\) 在分布意义（distributional sense）下满足拉普拉斯方程。换言之，对于任意测试函数 \(\varphi \in C_c^\infty(\Omega)\)，有\(\int_{\Omega} f \, \Delta \varphi \, dx = 0.\)根据魏尔引理（Weyl’s lemma），弱调和函数在几乎处处意义下与强调和函数（strongly harmonic function）一致，因此它是平滑的（即 \(C^\infty\)）。同样地，任何弱调和分布实际上都对应于某个强调和函数，因此也必然光滑。

狄利克雷原理（Dirichlet’s Principle）拉普拉斯方程还有若干其它等价的弱形式，其中之一即为著名的\emph{狄利克雷原理}。在该原理中，调和函数可视为属于 Sobolev 空间 \(H^1(\Omega)\) 的函数，它们是下述狄利克雷能量泛函的极小化者：
\[
J(u) := \int_{\Omega} |\nabla u|^2 \, dx.~
\]
换句话说，若\(u \in H^1(\Omega)\)满足\(J(u) \le J(u + v)\),对所有\(v \in C_c^\infty(\Omega)\),亦即对所有\(v \in H_0^1(\Omega)\)成立，则称 \(u\) 为调和函数。因此，调和函数可理解为使狄利克雷能量泛函在局部变分下取极小值的函数，即物理意义上的“平衡态势能场”。这一弱形式在变分法、偏微分方程以及数值分析中具有极其重要的作用。
\subsubsection{流形上的调和函数}
在任意黎曼流形（Riemannian manifold）上，调和函数可通过拉普拉斯–贝尔特拉米算子（Laplace–Beltrami operator）\(\Delta\) 来定义。在这种情形下，函数 \(f\) 称为调和的（harmonic），若满足
\[
\Delta f = 0.~
\]
许多欧氏空间中调和函数的性质在这种更一般的流形背景下依然成立，例如：平均值定理（Mean value theorem，基于测地球 geodesic balls）；最大值原理（Maximum principle）；哈纳克不等式（Harnack inequality）。除平均值定理外，这些性质均可视为一般线性二阶椭圆型偏微分方程的直接推论。
\subsubsection{次调和函数}
若一个 \(C^2\) 函数满足\(\Delta f \ge 0\),则称其为次调和函数（subharmonic function）。此条件保证了最大值原理依然成立，尽管其他调和函数的性质可能不再具备。更一般地，函数 \(f\) 为次调和，当且仅当：在其定义域内任意一个球的内部，其图像位于插值该球边界值的调和函数图像之下。
\subsubsection{调和形式}
调和函数研究的一个重要推广是调和形式（harmonic forms）的研究，它定义在黎曼流形上，并与上同调理论（cohomology theory）密切相关。此外，还可以定义向量值调和函数或调和映射（harmonic maps），即两个黎曼流形之间的映射，使其成为广义狄利克雷能量泛函\(E(u) = \int_M |\nabla u|^2 \, dV\)的临界点。这类映射包含普通的调和函数作为特例（即狄利克雷原理Dirichlet principle）。调和映射在极小曲面理论（theory of minimal surfaces）中尤为重要。例如，一条曲线（即从区间 \(I \subset \mathbb{R}\) 到黎曼流形 \(M\) 的映射）\(\gamma: I \to M\)是调和映射当且仅当它是一条测地线（geodesic）。








